\section{Introduction and main result}\label{sec:introduction}

For a Riemannian metric $g$, let $\Delta_g=\operatorname{div}_g\nabla_g$,
and write $Z_u=\{x:u(x)=0\}$ for the nodal set of a real function $u$.
We prove the following bound for one arbitrary fixed smooth surface.

\begin{theorem}\label{thm:main}
Let $(M,g)$ be a closed connected $C^\infty$ Riemannian surface.
There is a finite constant $C=C(M,g)$ such that every nonzero real
Laplace eigenfunction satisfying
\[
 -\Delta_g u=\lambda u,\qquad \lambda>0,
\]
obeys
\[
 \mathcal H_g^1(Z_u)\le C\sqrt\lambda.
\]
Here $\mathcal H_g^1$ is Hausdorff measure for the Riemannian distance.
The same constant applies to every eigenfunction in every positive
eigenspace.
\end{theorem}

Yau's conjecture asks for upper and lower bounds of order
$\sqrt\lambda$ for the $(n-1)$-dimensional measure of the nodal set
on a smooth closed $n$-dimensional manifold \cite{Yau1982}.
Donnelly and Fefferman proved both bounds for real analytic manifolds
and metrics \cite[Theorem~1.2]{DonnellyFefferman1988}.
For smooth surfaces, Br\"uning established the lower bound
\cite[Satz~2]{Bruning1978}; Logunov subsequently proved the lower bound
in every dimension \cite[Theorem~1.1]{Logunov2018Lower}.
For the surface upper bound with a smooth metric, Donnelly and
Fefferman proved $C\lambda^{3/4}$ \cite{DonnellyFefferman1990},
and Dong gave a different proof at the same order \cite{Dong1992}.
Logunov and Malinnikova then obtained
$C\lambda^{3/4-\beta}$ for a universal $\beta\in(0,1/4)$
\cite[Equation~(1.1)]{LogunovMalinnikova2018}.
Together with the known lower bound, Theorem~\ref{thm:main}
resolves Yau's conjecture positively for smooth closed surfaces.

\subsection*{Mean growth as the target}

Set $k=\sqrt\lambda$, the frequency. In local smooth isothermal
coordinates the metric is $g=p(x)(dx_1^2+dx_2^2)$, with $p>0$ smooth,
as in \cite[Section~1]{DeTurckKazdan1981}, and the eigenfunction
equation becomes
\[
 \Delta u+k^2p(x)u=0.
\]
Nazarov, Polterovich, and Sodin advocated a statistical approach to
local doubling and sketched its relation to nodal length, building on
Nadirashvili's approach
\cite[Section~7.3]{NazarovPolterovichSodin2005}.
Roy-Fortin established precise comparisons between growth and nodal
length \cite{RoyFortin2015,RoyFortinLq}. We prove the uniform mean
growth bound needed for the sharp upper estimate and use his planar
theorem for the final transfer.

Here is why an average suffices. For a coordinate square $Q$ of side
$s$ and center $x_Q$, consider the unweighted logarithmic growth
\[
 E(Q,0)=\log
 \frac{\|u\|_{L^2(\D_{Rs}(x_Q))}}
      {\|u\|_{L^2(\D_s(x_Q))}},
\]
where $R$ is the fixed large disk ratio chosen below and $\D_r(x)$ is
the Euclidean disk of radius $r$. The norms are unnormalized.
Subdivide a fixed coordinate square into congruent squares of side
$s$ comparable to $k^{-1}$, stopping once $ks$ is sufficiently small;
call these the terminal squares.
The estimate proved below gives
\[
 \av_{Q\text{ terminal}} E(Q,0)\le C.
\]
On each such square, Roy-Fortin's theorem, after rescaling and a
change from supremum norms to $L^2$ norms, gives
\[
 \mathcal H^1(Q\cap Z_u)\le Cs\bigl(1+E(Q,0)\bigr).
\]
There are $O(s^{-2})$ squares in a fixed coordinate region.
Summing the local estimates therefore gives $O(s^{-1})=O(k)$
nodal length. The precise grids, radii and constants are specified
in Sections~\ref{sec:subdivision} and~\ref{sec:length}.
The task is thus to control the average growth while allowing large
values on individual squares.

\subsection*{Affine correction and persistent growth}

The analytic step extracts logarithmic profiles from norms measured
on an exponential scale. With the required separation from the rescaled
frequency $ks$, a localized Carleman estimate makes these profiles
subharmonic, even after arbitrarily large affine tilts. Its decisive
feature is two-dimensional: when the normalized derivative energy
is perpendicular to the weight gradient, the commutator sees the
full trace of the weight Hessian.

A local affine part of a profile presents an obstruction to decay:
its oscillation and $ks$ decrease by the same factor
under subdivision. We subtract that affine part and carry its slope
separately. The remaining profile growth decreases faster than the
side length on average, while the changes of slope have bounded
average cost. Section~\ref{sec:subdivision} quantifies both estimates;
Section~\ref{sec:descent} sums their costs along paths that may stop.

To use this improvement through many scales, we divide a grid path
into excursions of high growth. A newly entered excursion produces
a nearby set of positive area on which high growth persists to the
terminal scale. These witness sets may lie outside their entry
squares. A packing argument controls them by counting discrepancies
between spatial membership and grid ancestry; those discrepancies
are confined to summable boundary collars. This geometric argument
is stated separately in Lemma~\ref{lem:packing}. It permits the
analytic persistence statement to be used without requiring its
witnesses to lie inside the original squares.

Section~\ref{sec:profiles} constructs the subharmonic limits and
establishes the initial growth bound. Section~\ref{sec:subdivision}
proves the affine subdivision estimate, and
Section~\ref{sec:descent} controls stopped descents and persistent
slopes. Section~\ref{sec:persistence} proves persistence at entry,
packs the excursions, and bounds the terminal mean growth.
Section~\ref{sec:length} completes the nodal-length estimate.
All coordinate constructions are local, so no orientability
assumption is needed.
